\selectlanguage{ngerman}%
\frenchspacing%
\begin{ledgroupsized}[r]{120mm}%
\footnotesize%
\pstart%
\noindent\textbf{Überlieferung:}%
\pend%
\end{ledgroupsized}%
\begin{ledgroupsized}[r]{114mm}%
\footnotesize%
\pstart%
\parindent -6mm%
\makebox[6mm][l]{\textit{L}}%
Aufzeichnung:
LH~XXXVII~5 Bl.\;124\textendash125.
Ein Bogen~4\textsuperscript{o},
oberer Rand beschnitten;
Papiererhaltungsmaßnahmen.
\href{\Scan{LH_37_05_124-125wz}}{Wasserzeichen} im Falz. %
Dreieinfünftel vollseitig beschriebene Seiten;
Bl.\;\href{\Scan{LH_37_05_125v}}{125\,v\textsuperscript{o}} ist um vier Fünftel unbeschrieben.
\pend%
\end{ledgroupsized}%
\vspace{5mm}%
\begin{ledgroup}%
\footnotesize%
\pstart%
\noindent%
\textbf{Datierungsgründe:}
Die Datierung der Aufzeichnung N.\;\ref{60339} beruht auf Eigenschaften des Textträgers sowie auf inhaltlichen Betrachtungen.
Siehe die ausführliche Begründung auf S.\;\pageref{60339_dat}.%
\pend%
\end{ledgroup}%
\vspace{8mm}
\selectlanguage{latin}%
\frenchspacing%
\normalsize%
\pstart%
\noindent%
\lbrack%
\href{\Scan{LH_37_05_124r}}%
{124\,r\textsuperscript{o}}%
\rbrack%
\pend %
\vspace*{-0.50em}%
\centerline{%
\includegraphics[width=0.19\textwidth]{%
gesamttex/edit_VIII-4/images/LH_37_05_124-125_d1_124r.pdf%
}}
\vspace{0.25em}
\centerline{%
\lbrack\textit{Fig.\:1}\rbrack%
}
\vspace{1.5em}%
\pstart%
\noindent%
\edlabel{LH_37_05_124-125_124r_methodus-1}%
Videor mihi
eccentrici corporum concursus%
\indexS{concursus eccentricus}%
\indexS{concursus corporum}
arcanum detexisse.%
\indexS{arcanum detectum}
\edtext{%
Ponamus concurrere
corpora \textit{A} et \textit{BC}%
\indexS{corpora concurrentia}
et lineam incursus perpendiculariter sumtam%
\indexS{linea incursus perpendicularis}
ad superficiem ipsius \textit{BC}%
\indexS{superficies corporis incurrentis}
esse \textit{AD}.
Jam a \textit{BC} abscindatur pars \textit{BE},%
}{%
\lemma{Ponamus}\Bfootnote{%
\textit{(1)}~corpus \textit{A} incurrere in corpus \textit{BC}
\textit{(2)}~concurrere corpora \textit{A} et \textit{BC},
\textit{(a)}~in d
\textit{(b)}~abscissamus\lbrack\textit{!}\rbrack\ a \textit{BC} partem
\textit{(c)}~et lineam %
\lbrack...\rbrack\ pars \textit{BE},%
~\textit{L}}}
cujus scilicet centrum gravitatis%
\indexS{centrum gravitatis partis abscissae}
\textit{D} incidit in rectam
\edtext{\textit{AD}.
His positis considero}{%
\lemma{\textit{AD}}\Bfootnote{%
\textit{(1)}~si opus
\textit{(2)}~. His positis
\textit{(a)}~concipio
\textit{(b)}~considero%
~\textit{L}}}
initio solum corpus \textit{BE},
et quis futurus sit eventus%
\indexS{eventus futurus}
si \textit{EC} non
\edtext{adesset.
Deinde fingo}{%
\lemma{adesset.}\Bfootnote{%
\textit{(1)}~Atque ita habebimus \textlangle mo\textrangle\
\textit{(2)}~Deinde fingo%
~\textit{L}}}
statim post nisum a concursu receptum%
\indexS{nisus receptus}%
\lbrack,\rbrack\
corpus \textit{BE} annexum intelligi%
\indexS{corpus annexum}
\edtext{ipsi \textit{EC}.%
\edlabel{LH_37_05_124-125_124r_centrgyr-1}
Vel}{%
\lemma{ipsi}\Bfootnote{%
\textit{EC}.
\textit{(1)}~Unde
\textit{(2)}~Vel%
~\textit{L}}}
si fingamus
\textit{BE} colligi in centr. gr.\;\textit{D},%
\indexS{corpus collectum in centro gravitatis}
et \textit{EC} in centr. gr.\;\textit{F},%
\indexS{corpus collectum in centro gravitatis}
perinde est %
ac si \textit{D} et \textit{F} %
connexa rigida \textit{DF},%
\indexS{linea rigida}
circa commune centrum gravitatis~\textit{H}%
\indexS{centrum gravitatis commune}
(\protect\vphantom)%
quod est centrum etiam totius \textit{BC}%
\protect\vphantom()
gyrari deberent,%
\indexS{gyratio circa centrum gravitatis}
ita tamen %
ut nonnihil etiam progrediatur totum
seu ipsum \textit{H}
\edlabel{LH_37_05_124r_ofgz-1}%
centrum.%
\edlabel{LH_37_05_124-125_124r_centrgyr-2}%
\indexS{centrum gravitatis commune}%
\edtext{}{%
\xxref{LH_37_05_124r_ofgz-1}{LH_37_05_124r_ofgz-2}%
\lemma{centrum.}\Bfootnote{%
\textit{(1)}~Hoc autem ex
\textit{(a)}~habitis
\textit{(b)}~datis
\textit{(2)}~Quomodo autem%
~\textit{L}}}
Quomodo autem%
\edlabel{LH_37_05_124r_ofgz-2}
nisus receptus%
\indexS{nisus receptus}
partim in motum totius,%
\indexS{motus totius}
partim in gyrationem%
\indexS{gyratio circa centrum gravitatis}
circa \textit{H} distribuatur,%
\indexS{nisus distributus}
sic definiri potest.
Nempe habetur jam motus ipsius \textit{A},
unde etiam regula habetur%
\indexS{regula reflexionis}%
\indexS{regula progressionis} %
quandonam reflectatur aut pergat,
quando scilicet reflecteretur aut pergeret %
si adesset solum \textit{BE}.
\ledpb%
Habetur
\edtext{%
praeterea summa virium absolutarum,%
\indexS{summa virium absolutarum}%
\indexS{vis absoluta}
eadem \lbrack scilicet\rbrack\
quae ante concursum.%
}{%
\lemma{praeterea}\Bfootnote{%
\textit{(1)}~motus
\textit{(2)}~summa virium
\bars{absolutarum \textit{erg.}} eadem
\bars{scilicet \textit{erg. Hrsg.}} quae ante concursum.%
~\textit{L}}}
Habetur et summa virium respectivarum,%
\indexS{summa virium respectivarum}%
\indexS{vis respectiva}
ea scilicet
quae foret
si solum concurrissent \textit{A} et
\edtext{\lbrack\textit{BE}\rbrack.}{%
\lemma{\textit{B}}\Bfootnote{%
\textit{L~ändert Hrsg.}}}
Haec autem nunc sic distribuenda est,%
\indexS{vis respectiva distribuenda}
ne augeatur postquam etiam \textit{EC} in motum%
\indexS{concitatio in motum}
\edtext{concitatur%
\lbrack;\rbrack\
nempe ipsorum \textit{D} et \textit{F},
seu \textit{BE} et \textit{EC}, gyratio
producit novam vim respectivam.%
\indexS{vis respectiva producta}%
\indexS{vis respectiva nova}%
}{%
\lemma{concitatur,}\Bfootnote{%
nempe
\textit{(1)}~concursus
\textit{(2)}~ipsorum
\textit{(a)}~gyratio
\textit{(b)}~\textit{D} et \textit{F} %
\lbrack...\rbrack\ gyratio producit
\textit{(aa)}~novum motum respectivum
\textit{(bb)}~novam vim respectivam.%
~\textit{L}}}
Seu concurrendo
\edtext{inter se \textit{D} et \textit{F}%
}{%
\lemma{inter}\Bfootnote{%
se \textit{D} et \textit{F}
\textit{erg.\;L}}}
possunt tendere determinatum
\edtext{Elastrum,%
\indexS{elastrum tensum}%
\indexS{tensio elastri}
si nempe fingamus a circulatione%
\indexS{circulatio in concursu eccentrico}
liberari %
et sic dirigi,
ut sibi directe occurrant.%
}{%
\lemma{Elastrum,}\Bfootnote{%
\textit{(1)}~et similiter concurrendo d
\textit{(2)}~(\protect\vphantom)si fingamus a
\textit{(3)}~ambo simul concurrendo
\textit{(4)}~si nempe %
\lbrack...\rbrack\ directe occurrant.%
~\textit{L}}}
Tenso hoc elastro%
\indexS{elastrum tensum}
coalescere ponantur%
\indexS{corpora coalescentia}
et concurrere
residua vi respectiva%
\indexS{vis respectiva residua}
ipsi tertio \textit{A},
\edtext{vi\lbrack re\rbrack s}{%
\lemma{vis}\Bfootnote{%
\textit{L~ändert Hrsg.}}}
autem illae respectivae simul,
una cum progressiva%
\indexS{vis progressiva}
quae est via centri grav. omnium,%
\indexS{centrum gravitatis omnium}
\edtext{dab\lbrack un\rbrack t}{%
\lemma{dabit}\Bfootnote{%
\textit{L~ändert Hrsg.}}}
vim totalem.%
\indexS{vis totalis}%
\edlabel{LH_37_05_124-125_124r_methodus-2}
\lbrack%
\href{\Scan{LH_37_05_124v}}%
{124\,v\textsuperscript{o}}%
\rbrack%
\pend%
\pstart%
\edlabel{LH_37_05_124-125_124v_verallg-1}
Videndum an haec meditandi ratio,%
\indexS{ratio meditandi}
ut scilicet exiguum temporis intervallum%
\indexS{intervallum temporis exiguum}
pro simultaneo substituamus,%
\indexS{intervallum temporis simultaneum}
etiam inservire nobis possit
in concursu corporum plurium%
\indexS{concursus plurimus}%
\indexS{concursus corporum}
inter se,
fingendo primum concurrere duo,
mox intervallo quantumlibet exiguo%
\indexS{intervallum quantumlibet exiguum}
tertium supervenire:
Ponendo incurrens%
\indexS{corpus incurrens}
in duo successive pergere
post incursum in prius.%
\indexS{incursus plurimus}
Quod si etiam ponatur
post concursum in prius %
reflecti;
attamen respectu sequentis,
quasi celerius
\edtext{%
sequen\lbrack s\rbrack,%
}{%
\lemma{sequentis}\Bfootnote{\textit{L~ändert Hrsg.}}}
potest novus poni concursus.%
\indexS{concursus novus}%
\indexS{concursus plurimus}%
\indexS{concursus corporum}
Videndum etiam
an idem possit prodire
diversis modis assignando
prius et posterius.%
\indexS{prius et posterius}
Et quando duo excipientia%
\indexS{corpus excipiens}
se habent homoeoptote ad incurrens,%
\indexS{corpus incurrens}
videamus an ea praelatio%
\indexS{praelatio}
quam pro arbitrio%
\indexS{arbitrium}
demus alterutri
\edtext{\lbrack excipi\rbrack entium,}{%
\lemma{incurrentium}\Bfootnote{%
\textit{L~ändert Hrsg.}}}
nullam introducat %
inaequalitatem in ipsis;%
\indexS{inaequalitas introducta}
quae est
ut ita dicam
instantia crucis,%
\indexS{instantia crucis}%
\indexS{crux}
seu proprius ratiocinandi modus.%
\indexS{modus ratiocinandi}
Nam si hoc examen%
\indexS{examen}
non ferret haec inquisitio,%
\indexS{inquisitio de concursu plurimo}
non foret proba,%
\indexS{proba de concursu plurimo}
saltem pro concursu plurium corporum.%
\indexS{concursus plurimus}%
\indexS{concursus corporum}
Posset tamen manere proba%
\indexS{proba de concursu eccentrico}
pro concursu Eccentrico.%
\indexS{concursus eccentricus}
\edtext{Habetur et alia proba,%
\indexS{proba alia}
nempe diversae suppositiones%
\indexS{suppositio}
debent dare idem.%
\edlabel{LH_37_05_124-125_124v_verallg-2}%
}{%
\lemma{Habetur}\Bfootnote{%
et \lbrack...\rbrack\ dare idem.
\textit{erg.\;L}}}
\pend%
\pstart%
Sed altera methodus%
\indexS{methodus pro concursu multiplicato}
pro concursu%
\indexS{concursus multiplicatus}
\edtext{multiplicato,
dum}{%
\lemma{multiplicato,}\Bfootnote{%
\textit{(1)}~tum
\textit{(2)}~dum%
~\textit{L}}}
duo vel plura quaevis conjunguntur in unum,%
\indexS{corpora conjuncta in unum}%
\indexS{conjunctio corporum in unum}
et ponitur post concursum %
\edtext{eadem vis progressiva%
\indexS{vis progressiva}%
}{%
\lemma{eadem}\Bfootnote{%
\textit{(1)}~vis absoluta
\textit{(2)}~vis progressiva%
~\textit{L}}}
et respectiva%
\indexS{vis respectiva eadem ante et post concursum}
quae ante,
est indubitata.%
\indexS{methodus indubitata}
Tantum videndum
an sit sufficiens ad rem determinandam.%
\indexS{methodus sufficiens ad rem determinandam}%
\indexS{res determinanda}
Nempe assumatur incognita tanquam data,%
\indexS{incognita assumta tanquam data}
nempe motus
qui cuique concurrentium%
\indexS{motus corporum post concursum}
dandus post concursum.
Et hinc calculando
\edtext{vidend\lbrack um\rbrack}{%
\lemma{videndo}\Bfootnote{%
\textit{L~ändert Hrsg.}}}
an secundum quamvis combinationem satisfiat.%
\indexS{combinatio quaevis}
\lbrack%
\href{\Scan{LH_37_05_125r}}%
{125\,r\textsuperscript{o}}%
\rbrack%
\pend%
\newpage%
\vspace{1.5em}%
\centerline{%
\includegraphics[width=0.39\textwidth]{%
gesamttex/edit_VIII-4/images/LH_37_05_124-125_d2_125r.pdf%
}}
\vspace{0.25em}
\centerline{%
\lbrack\textit{Fig.\:2}\rbrack%
}
\vspace{1.5em}%
\pstart%
\edlabel{LH_37_05_124-125_125r_conatuseccentr-1}%
Ab inquisitione%
\indexS{inquisitio de concursu eccentrico}
de concursu Eccentrico,%
\indexS{concursus eccentricus}
seu nisu vivo,%
\indexS{nisus vivus}%
\indexS{inquisitio de nisu vivo}
multum differt inquisitio de conatu Eccentrico mortuo.%
\indexS{conatus mortuus}%
\indexS{conatus eccentricus mortuus}%
\indexS{inquisitio de conatu mortuo}
Ut si corpus \textit{GLM},
cujus centrum gravitatis%
\indexS{centrum gravitatis corporis incumbentis}
sit \textit{G},
incumbat simul vectibus \textit{HL} et \textit{KM}%
\indexS{vectis tangens superficiem}
tangentibus superficiem%
\indexS{superficies corporis incumbentis}
\edtext{%
corporis.
Ducatur \textit{LN}%
}{%
\lemma{corporis.}\Bfootnote{%
\textit{(1)}~Jungatur
\textit{(2)}~Ducatur \textit{LN}%
~\textit{L}}}
perpendicularis ad contactum,%
\indexS{contactus corporis incumbentis et sustinentis}
utique nisus est in hac \textit{LN},%
\indexS{nisus mortus}
sed quia
\edtext{is non}{%
\lemma{is}\Bfootnote{%
\bars{non \textit{gestr.}} non%
~\textit{L}}}
pervenit in centrum,
quaeritur quomodo
\edtext{ejus actio sit aestimanda.%
\indexS{actio aestimanda}%
}{%
\lemma{ejus}\Bfootnote{%
\textit{(1)}~facultas aestim
\textit{(2)}~actio sit aestimanda.%
~\textit{L}}}
Per centrum gr. \textit{G}%
\indexS{centrum gravitatis corporis incumbentis}
ducatur verticalis \textit{GR}
quae ipsam \textit{LN} secet in \textit{F}.
Deinde juncta \textit{LG} %
ducatur angulo \textit{LGE} recto,%
\indexS{angulus rectus}
recta \textit{GE} %
ipsi \textit{LN} occurrens in \textit{E}.
Quaeritur
\edtext{an actio potentiae%
\indexS{actio potentiae}
in \textit{L},}{%
\lemma{an}\Bfootnote{%
\textit{(1)}~nisus
\textit{(2)}~actio
\bars{ipsius \textit{gestr.}} potentiae in \textit{L},%
~\textit{L}}}
secundum lineam \textit{LN} %
facta,
sit ita aestimanda%
\indexS{actio aestimanda}
ac si esset tantum ad
\edtext{ipsam
\lbrack in \textit{G}\rbrack,
ut \textit{LG}%
}{\lemma{ipsam}\Bfootnote{%
\bars{in \textit{erg. Hrsg.}}
\bars{\textit{L} \textit{ändert Hrsg.}},
ut \textit{LG}%
~\textit{L}}}
ad \textit{LE}.
An vero pro puncto \textit{G} %
adhibendo punctum~\textit{F},
\edtext{retinebimus totam actionem%
\indexS{actio tota}
\lbrack in\rbrack\ \textit{L},
perinde%
}{%
\lemma{retinebimus}\Bfootnote{%
\textit{(1)}~totam po
\textit{(2)}~totam actionem
\bars{in \textit{erg. Hrsg.}}
\textit{L}, perinde%
~\textit{L}}}
ac si punctum \textit{G}
sic esset impulsum %
\edtext{%
ut \textit{F}.
\indexS{centrum gravitatis impulsum}
Et ita similiter
si ad contactum%
\indexS{contactus corporis incumbentis et sustinentis}
\textit{M}%
}{%
\lemma{ut}\Bfootnote{%
\textit{G}
\textit{(1)}~, quo facto, si
\textit{(2)}~. Et ita similiter si ad
\textit{(a)}~punctum contactus
\textit{(b)}~contactum \textit{M}%
~\textit{L}}}
ducatur normalis \textit{MP}
\edtext{occurrens ipsi \textit{GR} in \textit{S},}{%
\lemma{occurrens}\Bfootnote{%
ipsi \textit{GR} in \textit{S}
\textit{erg.\;L}}}
et jungatur \textit{MG},
et angulo \textit{MGQ}
\edtext{recto%
\indexS{angulus rectus}%
}{%
\lemma{recto}\Bfootnote{%
\textit{erg.\;L}}}
\edtext{ducatur recta \textit{GQ},%
}{%
\lemma{ducatur}\Bfootnote{%
\textit{(1)}~normalis
\textit{(2)}~recta
\textit{(a)}~\textit{MG},
\textit{(b)}~\textit{GQ},%
~\textit{L}}}
quaeritur similiter %
an consideranda sit actio in
\edtext{\lbrack\textit{M}\rbrack}{%
\lemma{\textit{G}}\Bfootnote{%
\textit{L~ändert Hrsg.}}}
quasi esset ad absolutam,%
\indexS{actio absoluta}
ut \textit{MG} ad \textit{MQ};
an vero consideranda sit
ut facta in \textit{S},
et punctum \textit{G}
eodem modo affectum%
\indexS{centrum gravitatis affectum}
seu impulsum%
\indexS{centrum gravitatis impulsum}
intelligatur ut
\edtext{\lbrack\textit{S}\rbrack.}{%
\lemma{\textit{G}}\Bfootnote{%
\textit{L~ändert Hrsg.}}}%
\edlabel{LH_37_05_124-125-perinde}
\pend%
\pstart%
Vel quod
\edtext{eodem redit,
quaeritur an fingere liceat}{%
\lemma{eodem}\Bfootnote{%
\textit{(1)}~redit si fing
\textit{(2)}~redit, quaeritur an fingere liceat%
~\textit{L}}}
quasi \textit{KM} recta tangat in pluribus punctis %
ut in \textit{M} et \textit{T},
seu quasi raderet aliquandiu %
tota linea
\textit{TM}.
Et sane si poneretur
\textit{MGT} esse triangulum aequilaterum,%
\indexS{triangulum aequilaterum}
ut \textit{MG} et \textit{TG} essent aequales,
\ledpb%
tunc perinde esset haud dubie
ac si directio fuisset in \textit{G}.%
\indexS{directio centri gravitatis}
Sed si aliter res se haberet,
deberet ex hac hypothesi%
\indexS{hypothesis non approbata}
res eodem redire.
Si vero non approbamus hanc hypothesin,%
\indexS{hypothesis non approbata}
nova et subtilis inquisitio orietur,%
\indexS{inquisitio nova}%
\indexS{inquisitio subtilis}
eo casu,
praesertim,
quo contactus sint infiniti,%
\indexS{contactus infiniti}%
\indexS{casus contactuum infinitorum}
seu quo
\edtext{%
vectis superficiem aliquandiu lambit.%
\indexS{vectis lambens superficiem}%
}{%
\lemma{vectis}\Bfootnote{%
\textit{(1)}~radi
\textit{(2)}~radit
\textit{(3)}~superficiem aliquandiu lambit.%
~\textit{L}}}
\pend%
\pstart%
Superest et etiam haec consideratio,%
\indexS{consideratio}
quod fieri potest,
in solido,%
\indexS{solidum}
ut \textit{LN} non occurrat ipsi rectae \textit{GR}.
An ergo tunc
\edtext{%
similiter
fingemus quasi planum lambens%
\indexS{planum lambens}%
\indexS{planum perpendiculare}
perpendiculare ad %
\lbrack\textit{LN}\rbrack.%
\edlabel{LH_37_05_124-125_125r_conatuseccentr-2}%
}{%
\lemma{similiter}\Bfootnote{%
\textit{(1)}~\bars{pro \textlangle l\textrangle\ \textit{streicht Hrsg.}}
\textit{(2)}~fingemus quasi planum
\textit{(a)}~tangens\indexS{planum tangens}
\textit{(b)}~lambens perpendiculare ad
\bars{\textit{LM}. \textit{ändert Hrsg.}}%
~\textit{L}}}
\lbrack%
\href{\Scan{LH_37_05_125v}}%
{125\,v\textsuperscript{o}}%
\rbrack%
\pend%
\pstart%
\edlabel{LH_37_05_124-125_ende-1}%
Videamus
quid proferant istae Methodi,%
\indexS{methodus proferens}
in casibus praesertim facilibus;%
\indexS{casus facilis}
et videamus
an ipsa Methodus nobis examinum occasionem praebeat.%
\edlabel{LH_37_05_124-125_ende-2}%
\indexS{examen}%
\indexS{occasio examinis}
Deinde veniamus ad lapidem Lydium,%
\indexS{lapis Lydius}
qui consistit in Methodo generali%
\indexS{methodus generalis}
maximi possibilis descensus.%
\indexS{descensus maximus possibilis}
Ut appareat quid sit statuendum.
\pend%
\vspace{6.0em}%
\pstart%
\noindent%
\centering%
*     *     *
\pend%
\newpage%
\selectlanguage{german}%
\frenchspacing%
\footnotesize%
\pstart%
\noindent%
\textls{Ausführliche Datierungsgründe:}%
\pagelabel{60339_dat}
In der titellosen und undatierten Aufzeichnung N.\;\ref{60339} wird zunächst die Methode besprochen, die %
gedient haben soll, ein wolbehütetes \glqq Geheimnis\grqq\ (\textit{arcanum}) zu erschließen, nämlich das kinematische Verhalten zweier Körper im Fall ihres direkten exzentrischen Stoßes (\refpassage{LH_37_05_124-125_124r_methodus-1}{LH_37_05_124-125_124r_methodus-2}).
Diese Methode wird sodann ansatzweise auf den komplexeren Fall eines mehrfachen Stoßes angewendet (\refpassage{LH_37_05_124-125_124v_verallg-1}{LH_37_05_124-125_124v_verallg-2}).
Schließlich wird ein analoger Sachverhalt im Bereich der Statik untersucht: Ein unregelmäßiger Körper lastet mit exzentrischen Momenten auf zwei schiefen Ebenen (\refpassage{LH_37_05_124-125_125r_conatuseccentr-1}{LH_37_05_124-125_125r_conatuseccentr-2}). 
Das im Texträger anzutreffende Wasserzeichen %
ist im Leibniz-Nachlass nach heutigem Kenntnisstand nur für den Zeitraum 1690 bis Februar 1695 belegt.
Leibniz kehrte von seiner Italienreise\indexO{Italien} erst gegen Mitte Juni 1690 nach Hannover\indexO{Hannover} zurück (\textit{Chronik}, S.\;103).\cite{01236}
Die Aufzeichnung N.\;\ref{60339} ist somit wohl im Zeitraum zwischen seiner Rückkehr aus Italien\indexO{Italien} und dem Jahr 1695 niedergeschrieben worden.
Der Text bietet allerdings Anhaltspunkte dar, die eine weitere Präzisierung der relativen Chronologie und der Datierung ermöglichen.
\pend%
\pstart%
Bemerkenswert ist zunächst, dass %
N.\;\ref{60339} die Frage nach dem Mittelpunkt der von einem exzentrischen Stoß herrührenden Drehung unvermittelt mit der These beantwortet, ein exzentrisch angestoßener Körper drehe sich um seinen eigenen Schwerpunkt (\refpassage{LH_37_05_124-125_124r_centrgyr-1}{LH_37_05_124-125_124r_centrgyr-2}).
Hieraus ist zu schließen, dass N.\;\ref{60339} nach dem auf April bis spätestens Oktober 1691 datierbaren Entwurf N.\;\ref{60300} entstanden ist, in dem Leibniz zwar dieselbe Ansicht äußert, zudem aber bemerkt, in früheren Texten (etwa N.\;\ref{60142+60297} und N.\;\ref{60299}) habe er anders gedacht.
Diese Schlussfolgerung wird dadurch bestätigt, dass N.\;\ref{60300} auf Papier der Klostermühle zu Thierhaupten\indexO{Thierhaupten bei Augsburg} bei Augsburg\indexO{Augsburg} abgefasst ist (siehe die Datierungsgründe, S.\;\pageref{LH_37_05_044_datierung}).
\pend%
\pstart%
Relevant ist ferner die zu Beginn von N.\;\ref{60339} (\refpassage{LH_37_05_124-125_124r_methodus-1}{LH_37_05_124-125_124r_methodus-2}) beschriebene Methode, auf der die Erörterung des exzentrischen Stoßes im Text beruht:
Es wird angenommen, dass der exzentrisch angestoßene Körper sich im Augenblick des Stoßes in zwei Teile auflöse, deren eine zentral angestoßen werde, die andere unberührt bleibe; beide Teile würden sich im folgenden Augenblick wieder zusammenschließen.
Damit ließen sich im ersten Augenblick die Gesetze des direkten zentralen Stoßes anwenden, um nach dem Zusammenschluss der Teile das kinematische Verhalten des exzentrisch angestoßenes Körpers zu untersuchen. %
Diese fiktionale Methode, die bereits in den 1680er Jahren ansatzweise auf den exzentrischen Stoß angewendet worden war (siehe \textit{LSB} VIII,\;3 N.\;63, S.\;676.4\textendash12;\cite{01399} N.\;64, S.\;680.7\textendash8),\cite{01400} greifen auch die Aufzeichnungen N.\;\ref{58940} (\refpassage{LH_35_11_14_005-006_005r_methodus-1}{LH_35_11_14_005-006_005r_methodus-2}) und N.\;\ref{58939} (\refpassage{LH_35_11_14_003-004_003r_methodus-1}{LH_35_11_14_003-004_003r_methodus-2}) auf, die Leibniz eigenhändig auf den 31.\;Oktober\,/\,10.\;November 1690 datiert hat.
Dies ist Grund für die Vermutung, dass N.\;\ref{60339} etwa zur gleichen Zeit wie N.\;\ref{58940} und N.\;\ref{58939} verfasst wurde, nämlich als eine von ihnen herrührende oder sie begleitende methodologische Überlegung.
\pend%
\pstart%
Der Ausblick am Ende von N.\;\ref{60339} (\refpassage{LH_37_05_124-125_ende-1}{LH_37_05_124-125_ende-2}), in dem eine Erforschung \glqq einfacher Fälle\grqq\ exzentrischen Stoßes angedeutet wird, lässt sich schließlich als skizzenhafte Ankündigung der Untersuchungen betrachten, die in den eigenhändig auf September 1691 datierten Entwürfen N.\;\ref{60106} und N.\;\ref{60107} erfolgten.
(Der Schluss von N.\;\ref{60339} könnte zudem auf ein Gedankenexperiment zum Gleichgewicht in einem Fluidum wie in N.\;\ref{60107}, \refpassage{LH_35_10_13_008_008r_gedankenexp-1}{LH_35_10_13_008_008r_gedankenexp-2}, anspielen.)
Folglich dürfte die Aufzeichnung N.\;\ref{60339} vor N.\;\ref{60106} und N.\;\ref{60107} und somit insgesamt zwischen dem 10.\;November 1690 und September 1691 entstanden sein.
Eine frühere (ab der zweiten Hälfte Juni 1690) oder spätere Datierung (bis mindestens 1695) ist nicht auszuschließen.
\pend%
\selectlanguage{latin}%
\frenchspacing%
\normalsize%
